\documentclass[14pt, a4paper]{extreport}

\usepackage{kpi-lab}
% Defined by subject
\newcommand{\CourseTitle}{Бізнес-аналіз в ІТ}
\newcommand{\Teacher}{ст. в. Міщенко Л. Д.}
\newcommand{\ReportYear}{2026}
% Defined by Student
\newcommand{\StudentName}{Ковальов Олександр}
\newcommand{\StudentGroup}{ІМ-51мн}
\newcommand{\Variant}{4}
\newcommand{\CourseNumber}{2}
% Defined by Submission
\newcommand{\Topic}{Stakeholder Map та аналіз контексту}
\newcommand{\LabNumber}{1}

\begin{document}
	\pagenumbering{gobble}
	
	\begin{titlepage}
		\begin{center}
			\normalsize
			Національний технічний університет України\\
			«Київський політехнічний інститут імені Ігоря Сікорського» \\[0.5em]
			Факультет інформатики та обчислювальної техніки\\
			Кафедра обчислювальної техніки
			
			\vfill
			
			{\large \textbf{ЗВІТ}\\[0.5em]}
			{з лабораторної роботи №\LabNumber \\}
			{з дисципліни <<\CourseTitle>>}
			
			\vspace{2em}
			
			{\textbf{<<\Topic>>}}
			
			\vspace{1em}
			
			{\textbf{Варіант №\Variant}}
			
			\vfill
			
			\begin{flushright}
				Виконав: \\ студент \CourseNumber{} курсу, групи \StudentGroup \\
				\StudentName \\[1em]
				Перевірила: \\ \Teacher
			\end{flushright}
			
			\vfill
			
			\begin{flushright}
				Дата здачі: \DTMtoday
			\end{flushright}
			
			\vspace{5em}
			КИЇВ -- \ReportYear
		\end{center}
	\end{titlepage}
	
	\clearpage
	\pagenumbering{arabic}
	\setcounter{page}{2}
	
	\subsection*{Завдання}
	
	Метою цієї роботи є визначення контексту та стейкхолдерів для обраного проєкту, а також опис проблем та бізнес-можливостей (problem statement), які вирішуються для кінцевого користувача. Крім того, необхідно ідентифікувати основні групи стейкхолдерів та побудувати на основі цих даних Stakeholder Map.
	
	\section*{Хід роботи.}
	
	\textbf{1. Визначення контексту та бізнес-ідеї}
	
	\textbf{Назва проєкту:} XaiVolt. 
	
	\textbf{Суть продукту:} Апаратно-програмний комплекс (розумний контролер + мобільний застосунок), який об'єднує домашню електромережу, портативні зарядні станції (EcoFlow, Bluetti тощо), сонячні панелі та загальну мережу. Система використовує AI для аналізу графіків відключень, прогнозів погоди (для сонячної генерації) та тарифних зон, щоб автоматично перемикати джерела живлення та оптимізувати зарядку акумуляторів.
	
	\textbf{Цільовий ринок:} Власники приватних будинків та квартир, а також малий бізнес (кав'ярні, невеликі офіси, магазини), які потребують безперебійного живлення та хочуть оптимізувати витрати на електроенергію.
	
	\textbf{2. Опис проблеми та бізнес-можливості (Problem Statement)}
	
	Головна проблема полягає у складності ручного управління енергоресурсами: користувачі змушені самостійно перемикати рубильники, підключати зарядні станції та постійно слідкувати за графіками відключень. Це не лише незручно, але й створює значний ризик людської помилки. Крім того, виникає проблема неефективного споживання та фінансових втрат, оскільки зарядка акумуляторів або використання потужних приладів часто відбувається у пікові години з найвищими тарифами, замість того, щоб використовувати дешевший нічний тариф або накопичену сонячну енергію. 
	
	Ситуацію суттєво ускладнюють раптові знеструмлення, адже офіційні графіки відключень не завжди точні. Через це користувачі часто залишаються з розрядженими станціями в найбільш невідповідний момент, втрачаючи дорогоцінний робочий час, побутовий комфорт або, у випадку малого бізнесу, своїх клієнтів.
	
	\textbf{3. Ідентифікація груп стейкхолдерів та Stakeholder Map}
	
	Усіх стейкхолдерів було розділено на чотири ключові категорії: Внутрішня команда (Internal Team), Ухвалювачі рішень та Інвестори (Decision Makers \& Investors), Користувачі та Експерти (Users \& Experts), а також Зовнішні фактори (External / Enablers). На основі оцінки їхнього впливу (Influence) та інтересу (Interest) до проєкту було побудовано відповідну матрицю стейкхолдерів.
	
	\image{12cm}{map}
	
	Крім скріншоту, наведеного вище, її можна \href{https://miro.com/app/board/uXjVHgPsEzs=/?share_link_id=876354643353}{знайти в Miro}.
	
	Згідно з побудованою матрицею, у квадранті тісної взаємодії (Actively Engaged), який характеризується високим впливом та високим інтересом, знаходяться інвестори, Product Manager, CEO та ключові B2B клієнти, оскільки саме вони є головною рушійною силою проєкту. Групу, яку необхідно підтримувати задоволеною (Keep Satisfied) через їхній високий вплив, але низький інтерес, складають державні регуляторні органи, постачальники електроенергії та виробники зарядних станцій. Від їхніх дій безпосередньо залежить легальність та технічна сумісність продукту на ринку. 
	
	Тих, кого потрібно постійно тримати поінформованими (Keep Informed), об'єднує низький вплив при високому інтересі. До цієї категорії належать команда розробки, електрики-інсталятори та кінцеві власники будинків чи квартир (B2C), адже вони зацікавлені в успішному запуску та розширенні функціоналу, хоча й не визначають глобальну стратегію. Нарешті, у зоні мінімального моніторингу (Monitor) з низьким впливом та інтересом опинилися провайдери хмарної інфраструктури, постачальники електронних компонентів та потенційні конкуренти.
	
	\subsection*{Висновок}
	
	Під час виконання роботи було успішно проаналізовано бізнес-контекст стартапу XaiVolt. Чітко сформульований Problem Statement дозволив зрозуміти головні болі цільової аудиторії (B2B та B2C сегменти), пов'язані з нестабільністю електромереж та фінансовими втратами від неефективного енергоспоживання. 
	
	Побудована Stakeholder Map (Power vs. Interest Grid) дала змогу візуалізувати розстановку сил навколо проєкту. Вона показує, що найвищий пріоритет у комунікації та управлінні очікуваннями слід надавати інвесторам та ключовим B2B клієнтам, водночас суворо дотримуючись вимог регуляторів та постачальників електроенергії (Keep Satisfied). Цей аналіз є міцним фундаментом для подальшого збору вимог та розробки функціоналу системи XaiVolt.
	
\end{document}